\title{xLSTM: Extended Long Short-Term Memory}

\begin{abstract}
The foundational elements of Long Short-Term Memory (LSTM)—specifically, the constant error carousel and gating mechanisms—were established in the 1990s and have since enabled LSTMs to play a pivotal role in numerous advancements within deep learning, notably as the basis for the earliest Large Language Models (LLMs). Despite their enduring utility, the emergence of Transformers, characterized by parallel self-attention mechanisms, has signaled a significant shift by surpassing LSTMs in large-scale applications. This work revisits a fundamental question: To what extent can language modeling be advanced by scaling LSTMs to the billion-parameter regime, utilizing contemporary techniques from the modern LLM landscape while addressing LSTM shortcomings? To this end, we propose exponential gating augmented with suitable normalization and stabilization strategies. Additionally, we reengineer the LSTM’s memory, resulting in: (i) sLSTM, which incorporates a scalar memory and scalar update, together with a novel form of memory mixing; (ii) mLSTM, a design that allows full parallelism through the use of matrix memory and a covariance-based update mechanism. Embedding these variants into residual block frameworks produces xLSTM blocks, which can be stacked via residual connections to form complete xLSTM architectures. These improvements—exponential gating and revised memory architectures—substantially enhance the effectiveness of xLSTM, achieving competitive or superior performance relative to leading Transformer and State Space Model approaches in both accuracy and scalability.\\
Code available at: \url{https://github.com/NX-AI/xlstm}
\end{abstract}

\section{Introduction}
\label{sec:intro}

The Long Short-Term Memory (LSTM) concepts~\citep{Hochreiter:91,Hochreiter:97nips,Hochreiter:97}, namely the constant error carousel and the use of gating mechanisms, were developed as a solution to the vanishing gradient issue encountered in recurrent neural networks~\citep{Hochreiter:91,Hochreiter:00book}:

\begin{align}
\label{eq:lstm_idea}
  c_t \ = \ f_t \ c_{t-1} \ + \ 
          i_t \  z_t \ , \quad  
          h_t \ = \ o_t \ \psi({c_t}) \ . 
\end{align}

The constant error carousel involves incrementally updating the cell state $c_{t-1}$ (green) by adding cell inputs $z_t$, with this process regulated by sigmoid gates (blue). The input gate $i_t$ and forget gate $f_t$ manage the manner in which this update occurs, whereas the output gate $o_t$ determines the cell’s output, specifically the hidden state $h_t$. The cell state undergoes normalization or squashing through $\psi$, after which the output gate produces the hidden state.

LSTMs have achieved notable success across a wide range of fields \citep{Hochreiter:01,Hochreiter:07,Schmidhuber:15}, dominating the landscape of text generation up until the emergence of Transformers in 2017~\citep{Vaswani:17}. Their utility has been extensively validated in many sequence-driven applications, including but not limited to text generation~\citep{Graves:13,Karpathy:15blog}, handwriting synthesis~\citep{Graves:13}, sequence-to-sequence translation tasks~\citep{Sutskever:14nips}, program evaluation~\citep{Zaremba:14arxiv}, image captioning~\citep{Karpathy:15,Hossain:19}, automatic code generation~\citep{Karpathy:15blog}, modeling rainfall-runoff processes~\citep{Kratzert:18,Kratzert:19}, and building hydrological systems for flood warning~\citep{Nearing:24}. In the realm of reinforcement learning, LSTMs have set performance benchmarks for sequential models as demonstrated by the AlphaStar model in StarCraft II~\citep{Vinyals:17short}, OpenAI Five for Dota 2~\citep{OpenAI:19}, and in magnetic controller frameworks for nuclear fusion~\citep{Degrave:22short}. Their prowess in capturing abstract representations is especially pronounced, with LSTMs demonstrating proficiency at extracting and retaining semantic content in their internal memory~\citep{Karpathy:15blog}, as evidenced by the discovery of neurons specialized for numeracy and syntax~\citep{Lakretz:19}, linguistics~\citep{Bau:19}, and sentiment analysis~\citep{Radford:17arxiv}. Despite the advent of newer architectures, LSTMs continue to see practical deployment in critical real-world settings~\citep{Degrave:22short,Nearing:24} and have proven to be enduringly reliable.

Although LSTMs have achieved notable progress, they are subject to three primary shortcomings: (i) Irreversible storage choices. This drawback is illustrated through the {\it Nearest Neighbor Search} task (refer to Appendix~\ref{sec:appNearestNeighborSearch} for further details): Given a reference vector, a model must sequentially traverse a sequence to identify the most similar vector, then output its corresponding value at the end. As depicted in the left panel of Figure~\ref{fig:lstmProblems}, the mean squared error for this task demonstrates that LSTMs have difficulty overwriting a previously stored value upon encountering a closer match later in the sequence. In contrast, our proposed xLSTM overcomes this by employing exponential gating. (ii) Constrained memory capacity, requiring all information to be distilled into scalar cell states. This challenge is illustrated via the {\it Rare Token Prediction} task. The right panel of Figure~\ref{fig:lstmProblems} displays token prediction perplexity across frequency bands for Wikitext-103~\citep{Merity:17}. LSTMs perform suboptimally on infrequent tokens as a result of their limited capacity. xLSTM addresses this issue by utilizing a matrix-based memory. (iii) Inherent sequentiality induced by hidden-to-hidden state dependencies, precluding efficient parallel processing due to the intertwined memory updates from one timestep to the next.

The shortcomings inherent to LSTM architectures have contributed to the rise of Transformers~\citep{Vaswani:17} within the realm of language modeling. This raises the question: if these constraints are addressed and LSTMs are scaled to the proportions of present-day Large Language Models, what level of performance could they attain in language modeling tasks?

\section{Extended Long Short-Term Memory}
\label{sec:xLSTM}

In order to address the shortcomings of the LSTM, the Extended Long Short-Term Memory (xLSTM) framework proposes two key alterations to the original LSTM paradigm outlined in Equation~\eqref{eq:lstm_idea}. These adjustments—specifically, the incorporation of exponential gating and innovative memory architectures—lead to the addition of two new variants within the LSTM lineage: (i) the sLSTM (detailed in Section~\ref{sec:sLSTM}), which utilizes a scalar memory and scalar update mechanism alongside memory mixing, and (ii) the mLSTM (described in Section~\ref{sec:mLSTM}), which features a matrix-based memory system and employs a covariance-based (outer product) update that enables complete parallelization. Both models leverage exponential gating to augment the traditional LSTM architecture. The mLSTM achieves parallel execution by eliminating memory mixing, specifically by forgoing hidden-hidden recurrent pathways. Both sLSTM and mLSTM support extensions to architectures with multiple memory cells; sLSTM, in particular, allows for inter-cell memory mixing. Moreover, sLSTM can incorporate multiple heads, where memory mixing occurs solely among cells within individual heads, not across heads themselves, presenting a novel strategy for memory mixing when combined with exponential gating. In contrast, for mLSTM, utilizing multiple heads is functionally identical to expanding the number of memory cells.

When these novel LSTM variants are incorporated within residual block modules, they form xLSTM blocks (refer to Section~\ref{sec:xLSTMblock}). Arranging these xLSTM blocks in a residual stacking fashion leads to what we call xLSTM architectures (see Section~\ref{sec:xLSTMarchitecure}). The overall structure and its constituent parts are depicted in Figure~\ref{fig:xlstm_sketch}.

\subsection{Review of the Long Short-Term Memory}
\label{sec:orgLSTM}

The foundational concept of LSTM~\citep{Hochreiter:91,Hochreiter:97nips,Hochreiter:97} centers on the scalar memory cell, which serves as both the processing and storage core. This architecture is specifically designed to prevent vanishing gradients~\citep{Hochreiter:91,Hochreiter:00book} by employing what is termed the constant error carousel—reflected in the cell state update mechanism. Within the memory cell, three gates regulate information flow: the input, output, and forget gates, the last of which was first introduced by \citet{Gers:00}. The LSTM memory cell’s update equations at timestep $t$ are given by:

\begin{align}
c_t \ &= \  f_t \ c_{t-1} \ + \ i_t \ z_t &  & &\text{cell state} \\
h_t  \ &= \ o_t \ \tilde{h}_t \ , & \tilde{h}_t \ &= \ \psi \left( c_t\right)
  &\text{hidden state} \\
z_t \ &= \ \varphi \left( \tilde{z}_t \right) \ , 
  &\tilde{z}_t \ &=  \ \mathbf{w}^\top_{z} \ \mathbf{x}_t \ + \
  r_{z}  h_{t-1} \ + \  b_{z} \ \
  &\text{cell input} \\
i_t \ &= \ \sigma \left( \tilde{i}_t  \right) \ , 
  &\tilde{i}_t \ &= \ \mathbf{w}^\top_{i} \ \mathbf{x}_t \ + \
  r_{i}  \ h_{t-1} \ + \  b_{i} \ \
  &\text{input gate} \\
f_t \ &= \ \sigma \left(  \tilde{f}_t \right) \ , 
  &\tilde{f}_t \ &= \ \mathbf{w}^\top_{f} \ \mathbf{x}_t  \ + \
  r_{f}  \ h_{t-1} \ + \  b_{f} \ \
  &\text{forget gate} \\
o_t \ &= \ \sigma \left( \tilde{o}_t \right) \ , 
  &\tilde{o}_t  \ &= \ \mathbf{w}^\top_{o} \ \mathbf{x}_t \ + \
  r_{o}  \ h_{t-1} \ + \  b_{o} \ \
  &\text{output gate} 
\end{align}

The input weight vectors $\mathbf{w}_{z}$, $\mathbf{w}_{i}$, $\mathbf{w}_{f}$, and $\mathbf{w}_{o}$ connect the input $\mathbf{x}_t$ to the cell input, as well as to the input, forget, and output gates, respectively. Similarly, $r_{z}$, $r_{i}$, $r_{f}$, and $r_{o}$ denote the recurrent weight parameters linking the previous hidden state $h_{t-1}$ to the cell input and each of the gates. The biases $b_{z}$, $b_{i}$, $b_{f}$, and $b_{o}$ are associated with these components accordingly. Activation functions $\varphi$ and $\psi$ govern the cell input and hidden state, with $\tanh$ commonly employed for both. In particular, $\psi$ acts to scale or constrain the values of the cell state to prevent them from diverging. All gate activations utilize the sigmoid function $\sigma \left( x \right)= 1/(1+\exp(-x))$. Later advancements introduced vector-valued memory cells $\mathbf{c}_t \in \mathbb{R}^{d}$, facilitating the incorporation of recurrent weight matrices $\mathbf{R} \in \mathbb{R}^{d \times d}$, which blend the outputs from different memory cells~\citep{Greff:15}; further technical details can be found in Appendix~\ref{sec:apporgLSTM}. Ablation experiments indicate that each part of the memory cell architecture is indispensable~\citep{Greff:15}.

\subsection{sLSTM}
\label{sec:sLSTM}

To enable LSTMs to modify their storage choices, we incorporate exponential gates (highlighted in red), normalization, and stabilization. Specifically, both the input and forget gates may employ exponential activation functions. For normalization, we propose a normalizer state, which accumulates the products of the input gate with all subsequent forget gates.

The scalar sLSTM forward pass is given by:

\begin{align}
\label{eq:slstmforward}
c_t \ &= \  f_t \ c_{t-1} \ + \ i_t \ z_t & & 
 &\text{cell state} \\
n_t \ &= \  f_t \ n_{t-1} \ + \ i_t  & & 
 &\text{normalizer state} \\
h_t \ &= \ o_t \ \tilde{h}_t \ ,
  &\tilde{h}_t \ &= \ c_t / n_t
&\text{hidden state} \\
z_t \ &= \ \varphi \left( \tilde{z}_t \right) \ , 
  &\tilde{z}_t \ &=  \ \mathbf{w}^\top_{z} \ \mathbf{x}_t \ + \
  r_{z}  h_{t-1} \ + \  b_{z}
  &\text{cell input} \\
i_t \ &= \ \!\exp \left( \tilde{i}_t  \right) \ , 
  &\tilde{i}_t \ &= \ \mathbf{w}^\top_{i} \ \mathbf{x}_t \ + \
  r_{i}  \ h_{t-1} \ + \  b_{i}
  &\text{input gate} \\
f_t \ &= \ \sigma \left(  \tilde{f}_t \right) \ \text{OR} \
\!\exp \left(  \tilde{f}_t \right) \ ,
  &\tilde{f}_t \ &= \ \mathbf{w}^\top_{f} \ \mathbf{x}_t  \ + \
  r_{f}  \ h_{t-1} \ + \  b_{f}
  &\text{forget gate} \\
o_t \ &= \ \sigma \left( \tilde{o}_t \right) \ , 
  &\tilde{o}_t  \ &= \ \mathbf{w}^\top_{o} \ \mathbf{x}_t \ + \
  r_{o}  \ h_{t-1} \ + \  b_{o}
  &\text{output gate} 
\end{align}

We adapt the conventional LSTM gating methods—including input- or hidden-state-dependent gating along with a bias component—for integration into these novel architectures. Given that exponential activations are prone to producing excessively large outputs that may result in overflow, we employ an extra stabilizing state, denoted as \mbox{$m_t$~\citep{Milakov:18arxiv}}, to enhance gate stability:

\begin{align}
\label{eq:slstmstabil}
m_t \ &= \ \max \left( \log ( f_t ) + m_{t-1} , \log ( i_t ) \right) &\text{stabilizer state} \\
i'_t \ &= \ \!\exp \left( \log \left ( i_t \right) - m_t \right) \ = \ \!\exp \left( \tilde{i}_t - m_t \right) \ \, 
  &\text{stabil. input gate} \\
f'_t \ &= \ \!\exp \left( \log \left( f_t \right) + m_{t-1} - m_t \right) \ \, 
  &\text{stabil. forget gate}
\end{align}

As demonstrated in Appendix~\ref{sec:appsLSTM}, substituting $f_t$ with $f'_t$ and $i_t$ with $i'_t$ during the forward pass leaves both the network's overall output and the gradients of the loss with respect to the parameters unaffected.

\paragraph{New Memory Mixing.} The sLSTM architecture is capable of accommodating several memory cells, analogous to the original LSTM structure (refer to Appendix~\ref{sec:appsLSTM}). The inclusion of multiple memory cells facilitates the blending of memories through recurrent pathways $\mathbf{R}_{\mathbf{z}}$, $\mathbf{R}_{\mathbf{i}}$, $\mathbf{R}_{\mathbf{f}}$, $\mathbf{R}_{\mathbf{o}}$ that originate from the hidden state $\mathbf{h}$ and are directed toward the memory cell input $\mathbf{z}$ as well as the gating units $\mathbf{i}$, $\mathbf{f}$, $\mathbf{o}$. A novel feature in this mechanism arises from the use of exponential gating in memory mixing. The redesigned sLSTM may feature several heads, enabling intra-head memory mixing while preventing mixing between distinct heads. This combination of multi-head structure and exponential gating in sLSTM introduces a distinctive and novel form of memory mixing.

\subsection{mLSTM}
\label{sec:mLSTM}

In order to increase the storage capacity of LSTM networks, we generalize the memory cell from a single scalar $c \in \mathbb{R}$ to a matrix representation $\mathbf{C} \in \mathbb{R}^{d \times d}$. Consequently, memory retrieval operations involve matrix multiplication. Specifically, at each time step $t$, we intend to encode a pair of vectors—a key $\mathbf{k}_t \in \mathbb{R}^d$ and a value $\mathbf{v}_t \in \mathbb{R}^d$—adopting nomenclature from the Transformer architecture. Subsequently, at a later time $t + \tau$, the stored value $\mathbf{v}_t$ can be accessed using a query vector $\mathbf{q}_{t+\tau} \in \mathbb{R}^d$. This framework aligns with the principles of Bidirectional Associative Memories (BAMs) \citep{Kohonen:72,Anderson:72,Nakano:72,Anderson:77}.

To encode a key-value association, we employ the covariance update rule~\citep{Sejnowski:77,Dayan:91}, which is given by

\begin{align}
\mathbf{C}_t \ &= \ \mathbf{C}_{t-1} \ + \ \mathbf{v}_{t} \ \mathbf{k}_t^\top \ .
\end{align}

We postulate that a layer normalization is applied prior to mapping the inputs to keys and values, ensuring that these vectors possess zero mean. The covariance update procedure is theoretically optimal~\citep{Dayan:91} for maximizing the separability of retrieved binary vectors, which is tantamount to achieving the highest possible signal-to-noise ratio. Limiting retrieval operations to only pairwise associations and allowing for quadratic computational complexity, as in attention mechanisms~\citep{Krotov:16,Krotov:17,Ramsauer:21}, enables even greater separability. This same covariance update mechanism aligns with the Fast Weight Programmers paradigm~\citep{Schmidhuber:92ncfastweights,Schlag:21}, which were subsequently modified to include a fixed decay factor for $\mathbf{C}_{t-1}$ and a constant scaling factor for 
$\mathbf{v}_{t} \mathbf{k}_t ^\top$~\citep{Ba:16}.
Following this conceptual lineage, we incorporate the covariance update approach within the LSTM architecture, interpreting the forget gate as the decay coefficient, the input gate as the learning coefficient, and assigning the output gate the role of modulating the recalled vector.

In this matrix memory framework, the normalizer state comprises a weighted aggregation of key vectors, with each key vector scaled by the corresponding input gate and the subsequent forget gates. As before, the normalizer state captures the cumulative influence of the gates. Given that the dot product between the query and the normalizer state may approach zero, we take its absolute value and enforce a minimum threshold (usually set at 1.0), following the method described in~\citep{Sun:23arxiv}. The mLSTM forward computation proceeds as follows:

\begin{align}
\label{eq:mlstm_recurrent_begin}
\mathbf{C}_t \ &= \  f_t \ \mathbf{C}_{t-1} \ + \ 
  i_t \ \mathbf{v}_t \ \mathbf{k}_t^\top & &  &\text{cell state} \\
\mathbf{n}_t \ &= \  f_t \ \mathbf{n}_{t-1} \ + \ 
  i_t \ \mathbf{k}_t & &  &\text{normalizer state} \\
\mathbf{h}_t  \ &= \ \mathbf{o}_t \ \odot \ \tilde{\mathbf{h}}_t
\ , \qquad \qquad 
\tilde{\mathbf{h}}_t \ = \   \mathbf{C}_t \mathbf{q}_t \ / \ 
  \max \left\{ |\mathbf{n}_t^\top \mathbf{q}_t|, 1 \right\} 
   &&&\text{hidden state} \\
\mathbf{q}_t \ &= \ \mathbf{W}_q \ \mathbf{x}_t \ + \ \mathbf{b}_q  & &  &\text{query input} \\
\mathbf{k}_t \ &= \ \frac{1}{\sqrt{d}} \mathbf{W}_k \ \mathbf{x}_t \ + \ \mathbf{b}_k  & &  &\text{key input} \\
\mathbf{v}_t \ &= \ \mathbf{W}_v \ \mathbf{x}_t \ + \ \mathbf{b}_v  & &  &\text{value input} \\
i_t \ &= \ \!\exp \!  \! \left( \tilde{i}_t  \right) \ , 
 \qquad \qquad \ \ \ \ \,
  \tilde{i}_t \ = \ \mathbf{w}^\top_{i} \ \mathbf{x}_t \ + \  b_{i} \ \ \ 
  &&&\text{input gate} \\
f_t \ &= \ \sigma \!  \left(  \tilde{f}_t \right) \ \text{OR} \ \!\exp \!  \! \left(  \tilde{f}_t \right) , \ \ \: \!
\tilde{f}_t \ = \ \mathbf{w}^\top_{f} \ \mathbf{x}_t  \ + \
  b_{f}
  &&& \text{forget gate} \\
\label{eq:mlstm_recurrent_end}
\mathbf{o}_t \ &= \ \sigma \left( \tilde{\mathbf{o}}_t \right) \ , \qquad \qquad \quad \ \ \ \,
  \tilde{\mathbf{o}}_t  \ = \ \mathbf{W}_{\mathbf{o}} \ \mathbf{x}_t \ + \
  \mathbf{b}_{\mathbf{o}}
  &&&\text{output gate} 
\end{align}

The architecture of mLSTM permits the use of several memory cells, analogous to the standard LSTM design. In mLSTM, employing multiple heads is functionally equivalent to having multiple cells due to the absence of any memory mixing mechanism. To ensure the stability of the exponential gates within mLSTM, we adopt the same stabilization strategies utilized in sLSTM (see Equation~\ref{eq:slstmstabil}). Owing to the lack of memory mixing in mLSTM, its recurrence relation admits a reformulation in parallel form. Further technical details can be found in Appendix~\ref{sec:appmLSTM}.

\subsection{xLSTM Architecture}
\label{sec:xLSTMarch}

\paragraph{xLSTM Blocks.}
\label{sec:xLSTMblock}
An xLSTM block aims to encode past information in a non-linear fashion within a high-dimensional representation, thereby enhancing the ability to distinguish among diverse histories or contexts. Effective separation of these histories is essential for accurate prediction of the subsequent sequence element, such as the next token. We draw on Cover’s Theorem~\citep{Cover:65}, which asserts that patterns embedded non-linearly in a space of higher dimensionality have a greater likelihood of being linearly separable than those in the original space. Our investigation focuses on two residual block configurations: (i) A residual block featuring post up-projection—similar to the approach used in Transformers—that first forms a non-linear summary of the prior information in the original space, subsequently applies a linear projection to a higher-dimensional space, introduces a non-linear activation, and finally maps the result back linearly to the initial dimensionality; this architecture is illustrated in the left panel of Figure~\ref{fig:Backbones}
and is detailed in the third column of Figure~\ref{fig:xlstm_sketch}. 
A comprehensive schematic is provided in Figure~\ref{fig:appsLSTM_detailed}
in the appendix. (ii) A residual block employing pre up-projection—analogous to State Space Models—where the mapping into the high-dimensional space occurs first via a linear transformation,

The model first generates a non-linear summary of prior information within a high-dimensional space, followed by a linear transformation back to the input space. For xLSTM blocks equipped with an sLSTM, the post up-projection block is applied. Conversely, for xLSTM blocks with an mLSTM, the pre up-projection block is chosen due to the increased memory capacity afforded by the high-dimensional space. Further clarification can be found by consulting the left panel of Figure~\ref{fig:Backbones}, the third column of Figure~\ref{fig:xlstm_sketch}, or Figure~\ref{fig:appsLSTM_detailed} in the appendix.

\paragraph{xLSTM Architecture. }
\label{sec:xLSTMarchitecure}
The xLSTM architecture is formed through the residual stacking of modular components~\citep{Srivastava:15,He:16}. Our design adopts the prevalent pre-LayerNorm~\citep{Ba:16b} residual backbone found in modern Large Language Models. Refer to the rightmost column of Figure~\ref{fig:xlstm_sketch} for a visual representation.

\subsection{Memory and Speed Considerations}

In contrast to Transformers, xLSTM networks maintain linear computational complexity and constant memory requirements as sequence length increases. The compressive nature of xLSTM memory makes these networks particularly advantageous for deployment in industrial contexts and edge computing scenarios.

The mLSTM's memory operates without additional parameters; however, it incurs significant computational cost due to its use of a $d \times d$ matrix for both storage and updating. In our approach, we balance the memory size with computational demands. Importantly, since these matrix operations are highly parallelizable on GPUs, their impact on actual runtime is minimal.

Although mLSTM supports parallelization similarly to FlashAttention~\citep{Dao:22,Dao:23} and GLA~\citep{Yang:23arxiv}, sLSTM cannot be parallelized because of its reliance on memory mixing through hidden-to-hidden interactions. Nevertheless, we have engineered an efficient CUDA version that incorporates register-level GPU memory optimizations, resulting in performance that is generally no more than twice as slow as mLSTM.

\section{Related Work}
\label{sec:related}

\paragraph{Linear Attention.} Various strategies have been introduced to address the quadratic time complexity with respect to context length inherent in the original Transformer architecture, thereby achieving attention mechanisms that scale linearly with sequence length. Synthesizer generates synthetic attention scores independently of token interactions~\citep{Tay:20arxiv}. Linformer achieves self-attention through a low-rank factorization and further applies linear approximation techniques~\citep{Wang:20arxiv}. Linear Transformer reformulates attention so it operates linearly in the sequence length~\citep{Katharopoulos:20}. Performer approximates the attention softmax function in linear time by employing positive orthogonal random features~\citep{Choromanski:21}. Furthermore, conventional attention has been substituted with efficient long-range convolutional operations in models such as Structured Global Convolution (SGConv)~\citep{Li:22} and Hyena Hierarchy~\citep{Poli:23}.

\paragraph{State Space Models.} In recent years, State Space Models (SSMs) have gained considerable attention due to their linear complexity with respect to sequence length and their competitive performance relative to Transformers. Initial developments in this area include the Structured State Space sequence model (S4)~\citep{Gu:21}. Subsequent models have been introduced, such as the Diagonal State Space (DSS) model~\citep{Gupta:22}, Gated State Space (GSS) models~\citep{Mehta:22}, S5 model~\citep{Smith:22}, Bidirectional Gated SSM (BiGS)~\citep{Wang:22}, H3 model~\citep{Fu:23}, and Mamba~\citep{Gu:24arxiv}.

\paragraph{Recurrent Neural Networks.} Due to their linear scaling with context length, Recurrent Neural Networks (RNNs) have been proposed as alternatives to Transformers and attention mechanisms. In particular, RNN architectures employing Deep Linear Recurrent Units (LRUs) have achieved notable performance on language modeling tasks~\citep{Orvieto:23, De:24arxiv}. Similarly, both Hierarchically Gated Linear RNN (HGRN)~\citep{Qin:23} and its successor HGRN2~\citep{Qin:24arxiv} have demonstrated encouraging results. Among RNN-based approaches, RWKV~\citep{Peng:23arxivshort,Peng:24arxivshort} is prominent for large language modeling and has proven to be a strong competitor to Transformer models.

\paragraph{Gating.}
Gating constitutes a fundamental concept originally introduced in LSTM architectures, and has since been revisited and adapted within numerous contemporary methods. The mechanism of gating appears in models such as HGRN~\citep{Qin:23}, HGRN2~\citep{Qin:24arxiv}, Gated Linear Attention (GLA)~\citep{Yang:23arxiv}, Gated State Space (GSS) models~\citep{Mehta:22}, Bidirectional Gated SSM (BiGS)~\citep{Wang:22}, Moving Average Equipped Gated Attention (MEGA)~\citep{Ma:22}, RWKV~\citep{Peng:23arxivshort}, and Mamba~\citep{Gu:24arxiv}.

\paragraph{Covariance Update Rule.} In order to improve storage capabilities, we augmented the mLSTM cell by integrating a matrix memory that utilizes a covariance update rule. Similar update strategies have been employed in Fast Weight Programmers \citep{Schmidhuber:92ncfastweights,Schlag:21}, RWKV-5 and RWKV-6~\citep{Peng:24arxivshort}, Retention~\citep{Sun:23arxiv}, Linear Transformer~\citep{Katharopoulos:20}, as well as HGRN2~\citep{Qin:24arxiv}.

\paragraph{Most Related.} The models most comparable to xLSTM from a conceptual standpoint include Retention~\citep{Sun:23arxiv}, RWKV~\citep{Peng:23arxivshort,Peng:24arxivshort}, and HGRN2~\citep{Qin:24arxiv}, as they all employ mechanisms such as matrix memory and/or gating. Nonetheless, unlike the proposed sLSTM, these frameworks lack support for memory mixing. The capability for memory mixing facilitates the resolution of state tracking tasks, making LSTMs strictly more expressive than both State Space Models (SSMs) and Transformers~\citep{Merrill:24,Deletang:23}. State tracking functionality is essential for operations such as code evaluation or following entities throughout extensive narratives.

\paragraph{Residually Stacking Architectures.} The xLSTM architectures, consistent with the vast majority of recent large-scale deep learning models, are assembled through the residual stacking of fundamental modules~\citep{Srivastava:15,He:16}.

This approach has paved the way for significant advancements in deep convolutional neural networks~\citep{He:16} and underlies the design of Transformers~\citep{Vaswani:17}. Transformers, in turn, are the foundational technology for contemporary Large Language Models (LLMs), including GPT-3~\citep{Brown:20short}, ChatGPT~\citep{Schulman:22short}, GPT-4~\citep{Achiam:23gpt4short}, Megatron-LM~\citep{Shoeybi:19}, Gopher~\citep{Rae:21short}, ERNIE 3.0 Titan~\citep{Wang:21arxivshort}, GLaM~\citep{Du:21glamshort}, Chinese M6~\citep{Lin:21}, the multilingual AlexaTM 20B~\citep{Soltan:22}, OPT~\citep{Zhang:22opt}, Chinchilla~\citep{Hoffmann:22short}, BLOOM~\citep{Scao:22arxivshort}, GLM-130B~\citep{Zeng:22arxivshort}, LaMDA~\citep{Thoppilan:22short}, PaLM~\citep{Chowdhery:22short}, Llama~\citep{Touvron:23arxiv}, and Gemini~\citep{GeminiTeam:23arxiv,Reid:24arxiv}.

\section{Experiments}
\label{sec:Exp}

In this section, we present an empirical assessment of xLSTM and benchmark it against established approaches, with particular emphasis on language modeling tasks. Section~\ref{sec:ExpSyntethic} explores xLSTM’s distinctive properties by applying it to a suite of synthetic benchmarks. 
Section~\ref{sec:ExpComparison} offers a comparative analysis of various modern language modeling techniques by measuring their validation set perplexity following training on 15B tokens sourced from SlimPajama~\citep{Soboleva:23}. Additionally, ablation studies are conducted on xLSTM using this dataset to dissect its architectural contributions. The scaling characteristics of all evaluated methods are then analyzed in line with the methodologies described in \citet{Kaplan:20} and \citet{Brown:20short}.
A comprehensive language modeling evaluation is presented in Section~\ref{sec:ExpLanguage}, where xLSTM and the strongest contenders identified in Section~\ref{sec:ExpComparison} are trained with an extended corpus of 300B tokens from SlimPajama~\citep{Soboleva:23}. This phase of the study begins by examining model generalization to extended context lengths, proceeds to assess validation perplexity and downstream task performance following the procedures of~\citep{Sutawika:24}, evaluates model robustness across 571 diverse text domains using the PALOMA language benchmark~\citep{Magnusson:23arxivshort}, and concludes by revisiting scaling trends under the scenario of a twentyfold increase in training data.

Throughout all experiments, we denote the configuration xLSTM[$a$:$b$] to represent an arrangement where the proportion of mLSTM-based to sLSTM-based xLSTM blocks is $a/b$. For instance, xLSTM[7:1] specifies a total of eight blocks, seven utilizing mLSTM and one utilizing sLSTM. With a typical setting of 48 total blocks, this corresponds to 42 mLSTM-based blocks and 6 sLSTM-based blocks. Additionally, across all experimental setups, we consistently employ pre up-projection blocks for mLSTM and post up-projection blocks for sLSTM.

\subsection{Synthetic Tasks and Long Range Arena}
\label{sec:ExpSyntethic}
We begin by evaluating the benefits of xLSTM's novel exponential gating with memory mixing on formal language tasks~\citep{Deletang:23}. Subsequently, we examine how xLSTM's updated matrix memory performs in the Multi-Query Associative Recall task~\citep{Arora:23arxiv}. Lastly, we measure xLSTM’s capacity for handling extended sequences using the Long Range Arena benchmark~\citep{Tay:21}.

\paragraph{Evaluation of xLSTM's Exponential Gating Mechanism with Memory Mixing.} We evaluate the recently introduced exponential gating with memory mixing in xLSTM, which is hypothesized to facilitate the resolution of state tracking tasks~\citep{Merrill:24, Merrill:23}. Our study extends the formal language benchmarks introduced by \citet{Deletang:23} to include training on variable-length sequences, thereby permitting investigation of length extrapolation. A comprehensive overview of these tasks along with additional results is provided in Appendix~\ref{sec:appExpSynthetic-formalLang}. 

We conduct a comparative analysis between xLSTM and alternative approaches such as Transformers, State Space Models, and conventional Recurrent Neural Networks. Performance is quantified by computing accuracy specifically on task-critical tokens, normalized on a scale from 0 (random performance) to 1 (perfect performance). The benchmarks consider two-block variants from each model: xLSTM[0:1] (utilizing only the sLSTM component), xLSTM[1:0] (utilizing only the mLSTM component), xLSTM[1:1], Llama, Mamba, RWKV, Retention, Hyena, standard LSTM, and LSTM implemented within Transformer blocks (LSTM (Block)). The results are depicted in Figure~\ref{fig:formal-main}. 

Notably, architectures such as Transformers and State Space Models that lack mechanisms for memory mixing and state tracking are unable to solve certain regular grammar tasks, for example, the parity task. This observation corroborates prior work, which has established that Transformers and State Space Models possess fundamental limitations in comparison to RNNs~\citep{Merrill:24, Merrill:23, Deletang:23}.

\paragraph{Test of xLSTM's Memory Capacities on Associative Recall Tasks.} This study evaluates the novel matrix memory in xLSTM by assessing its ability to retain information in the Multi-Query Associative Recall task~\citep{Arora:23arxiv}. In each test sequence, a set of key-value pairs is randomly selected from an extensive vocabulary; the model must remember these associations for subsequent retrieval. To increase the task's complexity beyond the original formulation, we raise the number of key-value pairs to as many as 256 and expand the context window to 2048 tokens. These modifications provide a comprehensive benchmark for assessing memory capacity across various model architectures. The evaluation compares two-block versions of Llama, Mamba, RWKV-5, RWKV-6, xLSTM[1:1], and xLSTM[1:0], using the accuracy of retrieving key-value pairs as the primary metric. Given the exponential memory capacity in the embedding size for Transformer-based models such as Llama~\citep{Ramsauer:21}, these models serve as the reference for optimal associative recall performance. The outcomes are depicted in Figure~\ref{fig:mqar-main}. Notably, xLSTM[1:1] achieves the highest recall performance among non-Transformer models, including in smaller configurations. The integration of the sLSTM block, contrary to potentially limiting memory, actually augments xLSTM's capacity, particularly in the most demanding case involving 256 key-value pairs. Further analysis, including extrapolation studies that suggest xLSTM can generalize to contexts surpassing those encountered during training, is available in Appendix~\ref{sec:appExpSynthetic-mqar}.

\paragraph{Test of xLSTM's Long Context Capabilities on Long Range Arena.} To evaluate xLSTM's effectiveness with lengthy sequences and extensive context windows, we benchmark various approaches on the Long Range Arena~\citep{Tay:21}. Across all evaluated tasks, xLSTM achieves robust results, indicating that its architecture is highly capable when addressing a broad range of long-context challenges.

Additional information is provided in Appendix~\ref{sec:appExpSynthetic-lra}.

\subsection{Method Comparison and Ablation Study}
\label{sec:ExpComparison}

In order to investigate the central question posed in our work—namely, the capabilities of our proposed LSTM variants at scale in the context of language modeling—we conduct experiments where xLSTMs, along with Transformers, State Space Models, and various other approaches, are trained using 15B tokens extracted from the SlimPajama dataset under a consistent auto-regressive framework. We evaluate the resulting models using the validation split and carry out ablation studies specifically for xLSTMs to analyze their components.

\paragraph{Comparative Analysis of xLSTM and Alternative Approaches.} For the purposes of comparison, models are trained using a dataset of 15B tokens from SlimPajama~\citep{Soboleva:23}. Model performance is assessed through perplexity measurements on the corresponding validation dataset. The methods under evaluation include: xLSTM (our proposed approach), GPT-3 (Transformer architecture)~\citep{Brown:20short}, Llama (Transformer)~\citep{Touvron:23arxiv}, H3 (SSM)~\citep{Fu:23}, Mamba (SSM)~\citep{Gu:24arxiv}, RWKV-4, RWKV-5, RWKV-6 (RNN-based models)~\citep{Peng:23arxivshort, Peng:24arxivshort}, GLA (a variant of linear Transformer)~\citep{Yang:23arxiv}, HGRN and HGRN2 (RNNs)~\citep{Qin:23, Qin:24arxiv}, RetNet (linear Transformer)~\citep{Sun:23arxiv}, Hyena (linear Transformer)~\citep{Poli:23}, along with xLSTM[1:0] and xLSTM[7:1] (as described in Section~\ref{sec:xlstm_blocks_notation}). 

Training utilized mixed-precision across the board; however, for models RWKV-5, RWKV-6, GLA, and HGRN2, this mixed-precision approach was implemented without employing the automated mixed-precision facility provided by PyTorch (refer to Appendix Section~\ref{sec:appGenTrainProc} for further details). 

The comparative frameworks are organized as follows: (a) Transformers, (b) State Space Models (SSMs), and (c) Recurrent Neural Networks (RNNs) together with linear Transformers, which are characterized by a linearized replacement for the standard Transformer attention. All models are configured to be parameter-matched to a GPT-3 variant containing 350M parameters—this entails an embedding dimension of 1024 and a stack of 24 residual layers. Of note, only GPT-3 utilizes shared embeddings for both token and output layers, which results in a reduced parameter count.

Results summarized in Table~\ref{tab:spaj15b_model_results} demonstrate that xLSTM achieves superior validation perplexity compared to existing alternatives. For a comprehensive breakdown, please see Appendix~\ref{sec:appExpComparison}. Scaling properties for this suite of experiments are visualized in Figure~\ref{fig:temp_sclaw15B}, and suggest that xLSTM is expected to retain its performance advantages as model sizes increase.

\paragraph{Ablation Studies.}
As shown in Table~\ref{tab:spaj15b_model_results} and Figure~\ref{fig:temp_sclaw15B}, xLSTM delivers strong performance on the language modeling task when trained with 15B tokens from the SlimPajama dataset. To analyze the impact of various architectural modifications leading from LSTM to xLSTM, we incrementally adjust a standard LSTM model toward the xLSTM configuration. The process begins with embedding LSTM layers within pre-LayerNorm residual architectures. The next step incorporates these layers into a post up-projection structure. The final stage involves the inclusion of exponential gating as well as matrix memory mechanisms.

Table~\ref{tab:ablstudies} (top) presents these findings.

The ablation experiments indicate that the enhanced performance stems from the synergy between exponential gating and the matrix memory. Furthermore, considering the critical role of gating in both RNNs and State Space Models, various gating architectures are analyzed. As seen in Table~\ref{tab:ablstudies} (bottom), making each gate trainable and responsive to the input leads to further performance gains. Further details regarding experiments on individual backbone modules can be found in Appendix~\ref{sec:appAblDetails}.

\subsection{xLSTM as Large Language Model}
\label{sec:ExpLanguage}

To conclude our investigation, we perform extensive language modeling experiments to evaluate xLSTM's viability as a large language model (LLM). For this purpose, we scale up the dataset, utilizing 300B tokens drawn from SlimPajama for training. This token count matches the data used by models such as Mamba~\citep{Gu:24arxiv} and Griffin~\citep{De:24arxiv}.

We conduct a comparison between xLSTM, RWKV-4, Llama, and Mamba, selecting one representative from each approach as described in Section~\ref{sec:ExpComparison}. RWKV-4 serves as the exemplar for RNN-based methods because viable training precision settings for RWKV-5, RWKV-6, and HGRN2 were established only after the initiation of experiments involving 300B tokens (refer to Appendix~\ref{sec:appGenTrainProc} for further details). Our experiments encompass multiple model sizes (125M, 350M, 760M, 1.3B), evaluating every model for their ability to extrapolate to longer sequences, and measuring their validation set performance. We further examine downstream task outcomes, assess language modeling results across 571 domains included in the PALOMA benchmark, and analyze model behavior under scaling laws.

\paragraph{Sequence Length Extrapolation.}
\label{sec:spaj300B_ctx_extrapolation}
To begin, we evaluate the ability of large-scale models—specifically, xLSTM, RWKV-4, Llama, and Mamba with 1.3B parameters—to generalize to longer sequences. Each model undergoes training on sequences of length 2048 and is subsequently evaluated on context lengths that extend up to 16384. Results are presented in Figure~\ref{fig:spaj300B_seq_extrapolation}. Notably, among these approaches, xLSTM consistently achieves lower perplexity as the context window grows.

\paragraph{Validation Perplexity and Downstream Tasks.}
\label{sec:spaj_downstream_eval}
Additionally, across every model scale, we assess xLSTM, RWKV-4, Llama, and Mamba models using the SlimPajama validation dataset to determine next token prediction abilities, as well as on downstream evaluations targeting common sense reasoning capabilities.

The third column in Table~\ref{tab:lmeval} presents the validation set perplexities corresponding to various approaches. Notably, both xLSTM[1:0] and xLSTM[7:1] achieve the lowest validation perplexities across all evaluated model sizes. The remaining columns in Table~\ref{tab:lmeval} summarize downstream task performance. Across nearly all tasks and model configurations, xLSTM consistently outperforms other methods—Mamba surpasses xLSTM solely on the ARC task in certain scenarios. Additional experimental details are available in Appendix~\ref{sec:appExpLanguage}.

\paragraph{Performance on PALOMA Language Tasks.} Third, we evaluate the ability of xLSTM, RWKV-4, Llama, and Mamba models of varying scales to predict the next token within PALOMA language tasks~\citep{Magnusson:23arxivshort}. Evaluation is conducted by calculating perplexity scores for next token prediction across 571 distinct text domains, encompassing sources such as nytimes.com and Reddit’s r/depression community. Table~\ref{tab:ppleval} presents results with perplexity scores separated into broad language modeling (the first seven columns) and specialized domain assessments (the last five columns). Across these language benchmarks, the xLSTM[1:0] configuration consistently outperforms the xLSTM[7:1] setup.

xLSTM[1:0] demonstrates lower perplexity compared to Mamba in 568 out of 571 (99.5\%) text domains, outperforms Llama in 486 of 571 (85.1\%) cases, and achieves lower perplexity than RWKV-4 in 570 out of 571 (99.8\%) instances. Further information can be found in Appendix~\ref{sec:appExpLanguage}.

\paragraph{Scaling Laws.} Fourth, we analyze the scaling properties governed by power-law relationships, which enable predictions of model performance at increased parameter counts~\citep{Kaplan:20,Brown:20short}. Figure~\ref{fig:temp_sclaw300B} displays the observed scaling trends. Although all architectures exhibit comparable scaling patterns, their baselines vary. RWKV-4 consistently yields the lowest performance, followed sequentially by Llama and Mamba. xLSTM outperforms Mamba by a margin similar to the gap between Mamba and Llama. The observed scaling suggests that xLSTM is likely to retain or extend its performance advantage over both Transformers and State-Space models as model size increases.

\textbf{Generation Times and Maximal Throughput.} In our evaluation, as depicted in Figure~\ref{fig:inference_time_generation}, we assess the text generation latency for various xLSTM architectures at the 1.3B parameter scale, and further analyze the maximum achievable throughput in Figure~\ref{fig:inference_time_decoding_throughput} (left). Our comparisons include similarly sized implementations of Mamba, Llama, and RWKV models from HuggingFace, with the Llama baseline utilizing a static key-value cache. It should be noted that, during our experimental phase, neither full cache compilation for Transformer-based models nor the use of \texttt{torch.compile} for Mamba could be executed successfully. All text generation benchmarks were conducted using a batch size of 1 and a pre-fill value of 16, conditions especially advantageous for Transformer architectures.

In Figure~\ref{fig:inference_time_generation}, it becomes apparent that xLSTM, alongside recurrent baselines Mamba and RWKV-4, exhibits a linear growth in generation time with sequence length, in contrast to the quadratic scaling observed in the Llama model. Regarding decoding throughput, we varied both batch sizes and pre-fill parameters specifically for the Llama configuration. As demonstrated in Figure~\ref{fig:inference_time_decoding_throughput} (right), xLSTM is capable of operating with considerably larger batch sizes than Llama due to its fixed memory usage, thereby achieving superior throughput levels.

\section{Limitations}

(i) Unlike mLSTM, sLSTM’s memory mixing mechanism precludes parallelized computation, thereby restricting the potential for highly efficient parallel implementations. However, we have engineered a CUDA kernel for sLSTM that currently achieves performance within a factor of two compared to our parallelized mLSTM kernel. (ii) The existing CUDA kernels for mLSTM lack optimization, resulting in a runtime approximately four times greater than that of FlashAttention or the scan operation in Mamba. Substantial speedups could be achieved by adopting optimization techniques similar to those in FlashAttention. (iii) In mLSTM, manipulating matrix memory entails processing $d \times d$ matrices, which is computationally intensive. Nevertheless, since both updating and accessing memory are parameter-free and compatible with parallel matrix computation, the additional wall clock time incurred by this complexity is relatively modest. (iv) Proper selection of the initial values for the forget gates is crucial. (v) Because the matrix memory size does not grow with the length of the input sequence, scaling up the sequence length could exceed the memory’s capacity when handling longer contexts. Despite this, no significant limitations have emerged for sequences up to 16k in length; see Section~\ref{sec:spaj300B_ctx_extrapolation}. (vi) Due to the prohibitive compute demands associated with large-scale language modeling, we did not undertake extensive tuning of the architecture or hyperparameters, particularly for more sizable xLSTM models. We expect that comprehensive architectural and hyperparameter optimization will be required for xLSTM to fully realize its capabilities.

\section{Conclusion}
We have provided a partial response to our central inquiry: To what extent can language modeling advance when LSTMs are scaled to the order of billions of parameters? Our findings suggest that the answer is: ``At minimum, on par with prevailing architectures such as Transformers and State Space Models.'' Our work introduces xLSTM, which extends LSTM through exponential gating mechanisms involving memory mixing and introduces a novel memory configuration. Empirically, xLSTM models yield competitive performance in language modeling tasks relative to established approaches including Transformers and State Space Models. Observed scaling trends imply that larger-scale xLSTM variants could emerge as formidable alternatives to current Transformer-based Large Language Models. Beyond language modeling, xLSTM exhibits promise for substantial influence in domains such as Reinforcement Learning, Time Series Prediction, and modeling of physical phenomena.

\end{document}